\documentclass[11pt]{article}
\usepackage[a4paper,margin=24mm]{geometry}
\usepackage{amsmath,amssymb,hyperref}
\title{Hugmap: contact branches and continuity}
\date{5 November 2026}
\begin{document}\maketitle
\section{What the function describes}
This is a static, planar circle-contact model, not a proof of convergence of a moving grasp.
Use a separate function for each fixed roll $r\in\{-90^\circ,+90^\circ\}$.
The physical logical bend satisfies $-\pi\le q_2<0$, and $|q_1|\le\pi$.
Changing roll is not a continuous operation within this planar model.
For input $z=(x,y,D)$ in millimetres, the intrinsic object centre and radius are
\[
 c_r(z)=\begin{pmatrix}(r/90^\circ)x/1000\\y/1000-0.1445\end{pmatrix},
 \qquad R=D/2000.
\]
The housing-centre origin is not the physical Go2 base origin.

\section{A finite composition of geometric functions}
Let $\ell=0.05$ metres and
\[
 S_n(a)=\ell\sum_{j=1}^{n}\begin{pmatrix}-\sin(ja)\\\cos(ja)\end{pmatrix}.
\]
Equivalently, away from removable singularities,
\[
 S_n(a)=\ell\frac{\sin(na/2)}{\sin(a/2)}
 \begin{pmatrix}-\sin((n+1)a/2)\\\cos((n+1)a/2)\end{pmatrix}.
\]
The five-link segment endpoints are $p_5=S_5(q_1/5)$ and
$p_{10}=p_5+\mathcal R(q_1)S_5(q_2/5)$.
Intermediate vertices and edge normals follow from these same rigid transforms
and the fixed CAD polygons. Thus no fitted spatial basis is needed.

For a chosen contact feature $a$, define its contact equation by
\[
 F_a(q,c,R)=
 \begin{cases}
 \|c-v_a(q)\|^2-R^2,&\text{vertex},\\
 n_a(q)^T(c-v_a(q))-R,&\text{edge with outward unit normal}.
 \end{cases}
\]
The finite edge interval, vertex normal cone, nonpenetration, protected-part
clearance, self-collision and joint limits are additional validity conditions.
For $a$ on segment 1 and $b$ on segment 2 or the gripper, solve
\[
 F_a(q_1,c_r(z),R)=0,\qquad F_b(q_1,q_2,c_r(z),R)=0.
\]
With $u=\tan(q_1/10)$ and $v=\tan(q_2/10)$, these become sequential polynomials
of degree at most ten. The implemented branch function is
\[
 Q_{r,a,b,j,k}(z)=
 \begin{pmatrix}
 10\arctan\bigl(u_j(z)\bigr)\\
 10\arctan\bigl(v_k(z,u_j(z))\bigr)
 \end{pmatrix}.
\]
Here $u_j$ and $v_k$ are admissible real roots. Root indices alone are not stable
branch identifiers; continuation must match joint configurations and contact domains.
The roots are evaluated numerically. This is a piecewise algebraic/trigonometric
construction, not a radicals-only closed form.

For two distinct circle contact points, write $d=n_1^Tn_2$, where the normals
are unit radial normals, and define
\[
 M(Q,z)=\sqrt{\frac{1+d}{1-d}}.
\]
This is the planar two-point friction threshold; antipodal contacts have zero
threshold, while coincident contact directions are singular. The implementation
checks the actual contact points and may choose the best pair when more than two
parts touch. The physical grasp still needs a strict friction margin.
The branch functions combine as
\[
 f_r(z)=\min_{Q\in\mathcal A_r(z)} M(Q,z),\qquad
 \mathcal B_r(\mu)=\{z:f_r(z)<\mu\},
\]
where $\mathcal A_r(z)$ is the set of valid enumerated contact roots. A missing
catalog root is reported as unknown coverage, not a proof of physical impossibility.

\section{Precisely when a branch is continuous}
Inside one contact domain, with strict collision and joint-limit margins,
\[
 J_q=
 \begin{pmatrix}
 \partial F_a/\partial q_1&0\\
 \partial F_b/\partial q_1&\partial F_b/\partial q_2
 \end{pmatrix}.
\]
If $\det J_q\ne0$, the implicit function theorem gives a locally smooth joint
branch $Q(z)$, with
\[
 \frac{\partial Q}{\partial z}=-J_q^{-1}\frac{\partial F}{\partial z}.
\]
This is a conditional local statement, not a certificate for the entire map.
There are several distinct events:
\begin{enumerate}
\item $M=0.5$: the posture can remain continuous while crossing out of the
allowed-friction set. A gap in the coloured map need not be a gap in geometry.
\item A feature changes from edge to vertex: two descriptions may join at the
same posture; the derivative can change without a posture jump.
\item Two candidate costs cross: the minimum cost can remain continuous while
selecting the cheapest posture causes a jump between distinct branches.
\item A collision, joint limit or singular root terminates a valid branch.
There may be no allowed continuation along the selected input path.
\end{enumerate}
The minimum of finitely many continuous costs is continuous where those branches
remain valid on a common neighbourhood. Their validity domains cannot be ignored.

\section{Numerical evidence and its limits}
The companion audit examines saved neighbours in $x,y,D$ separately at each roll,
then solves intermediate points rather than interpolating joint angles.
It computes the path minimizing the largest adjacent joint jump over the sampled
candidate layers, both with and without the friction cutoff. Refinement tests
whether the jump decreases with sampling distance.
This tests selected straight paths, not all paths in three-dimensional input space.
Neither the empirical feature catalog nor floating-point polynomial roots certify
exhaustive feasibility, and static continuation does not establish dynamic stability.
See \texttt{kyung-continuity/continuity.json} for measured results and the HTML
report for the posture and friction curves.
\end{document}
